% -*- coding: utf-8 ; -*-
\documentclass[dvipsnames]{article}% dvipsnames is for xcolor (loaded by Tikz, loaded by nicematrix)
\usepackage{xltxtra}
\usepackage[french]{babel}
\frenchsetup{og = « , fg = »}

\usepackage{xcolor}
\usepackage{tikz}
\usetikzlibrary{calc}
\usepackage{geometry}
\geometry{left=2.8cm,right=2.8cm,top=2.5cm,bottom=2.5cm,papersize={21cm,29.7cm}}

\usepackage{cascade}


\usepackage{enumitem}
\usepackage{verbatim}
\usepackage{amsmath}

% We use \MakeShortVerb of shortvrb and not \DefineShortVerb of fancyvrb
% because we don't want the contents of short verbatim colored in gray
\usepackage{shortvrb}
\MakeShortVerb{\|}

\setlength{\fboxsep}{0pt}

\usepackage{fancyvrb}
\fvset{commandchars=\~\#\@,formatcom=\color{gray}}

\usepackage{titlesec}
\titlespacing*{\section}{0pt}{6.5ex plus 1ex minus .2ex}{4.3ex plus .2ex}
\titlespacing*{\subsection}{0pt}{4.5ex plus 1ex minus .2ex}{2ex plus .2ex}


\def\interitem{\vspace{7mm plus 2 mm minus 3mm}}          
\def\emphase{\bgroup\color{RoyalPurple}\let\next=}

\usepackage{footnote}

\usepackage{varwidth}

\usepackage[hyperfootnotes = false]{hyperref}

\hypersetup
  {
    pdfinfo = 
      {
        Title = L’extension cascade ,
        Subject = Une extension LaTeX ,
        Author = F. Pantigny 
      }
  } 


\NewDocumentEnvironment {scope} {} {} {}

\NewDocumentCommand {\pkg} {m} {\textsf{#1}}
\NewDocumentCommand {\cls} {m} {\textsf{#1}}

\setlength{\parindent}{0pt}


\begin{document}

\VerbatimFootnotes


\title{L'extension LaTeX \pkg{cascade}\thanks{Ce document correspond à la 
version~\myfileversion\space de \pkg{cascade}, à la date du~\myfiledate.}} 
\author{F. Pantigny \\ \texttt{fpantigny@wanadoo.fr}}

\maketitle

\begin{abstract}
L'extension LaTeX \pkg{cascade} fournit une commande |\Cascade| pour faire des constructions qui présentent des
démonstrations mathématiques avec des accolades enchaînées pour les déductions. L'extension \pkg{cascade} fournit
également une commande |\Edacsac| (\emph{cascade} à l'envers) qui construit une structure similaire mais avec des
accolades allant dans l'autre sens. 
\end{abstract}



\vspace{1cm}
\section{La commande \textbackslash Cascade}

L'extension \pkg{cascade} fournit une commande |\Cascade| qui permet des
constructions comme la suivante, où la taille de l'accolade de droite est
calculée sur seulement une partie des éléments LaTeX composés à gauche.
\begin{center}
\Cascade{$\det(A) = \begin{vmatrix}3&4\\ -1&7\end{vmatrix}\neq 0$}
        {et, par conséquent, $A$ est inversible}
        {}
        {or $AX=Y$} 
donc, $X = A^{-1}Y$
\end{center}
        
\medskip
\begin{Verbatim}
~emphase#\Cascade@{$\det(A) = \begin{vmatrix}3&4\\ -1&7\end{vmatrix}\neq 0$}
        {et, par conséquent, $A$ est inversible}
        {}
        {or $AX=Y$} 
donc, $X = A^{-1}Y$
\end{Verbatim}


\interitem
La commande |\Cascade| prend ses quatre arguments comme suit :
\begin{center}
\Cascade{\fbox{\hbox to 2cm{\vrule height 18 pt depth 5 pt width 0pt\hfil \texttt{\#1}\hfil}}}%
{\fbox{\hbox to 3cm{\vrule height 10 pt depth 12 pt width 0pt\hfil \texttt{\#2}\hfil}}}%
{\fbox{\hbox to 2.5cm{\vrule height 15 pt depth 14 pt width 0pt\hfil \texttt{\#3}\hfil}}}%
{\fbox{\hbox to 1.7cm{\vrule height 12 pt depth 5 pt width 0pt\hfil \texttt{\#4}\hfil}}}
\end{center}

\interitem
Les commandes |\Cascade| peuvent être imbriquées comme dans l'exemple suivant :
\begin{center}
\Cascade{\ShortCascade{$(BH) \perp (AC)$}{$(OC) \perp (AC)$}}
        {or\enskip $(BH) \parallel (OC)$}
        {\ShortCascade{$(CH) \perp (AB)$}{$(OB) \perp (AB)$}}
        {or\enskip $(CH) \parallel (OB)$}
donc $(OBHC)$ est un parallélogramme
\end{center}

\interitem
Pour la lisibilité de ces constructions, une version simplifiée de |\Cascade|
est disponible, nommée |\ShortCascade|.  


Le code |\ShortCascade{X}{Y}| est simplement un raccourci pour le code
|\Cascade{}{X}{}{Y}|.

\medskip
L'exemple précédent peut être codé avec deux commandes |\ShortCascade|
et une commande |\Cascade| englobante.
\begin{Verbatim}
~emphase#\Cascade@{~emphase#\ShortCascade@{$(BH) \perp (AC)$}
                      {$(OC) \perp (AC)$}}
        {or\enskip $(BH) \parallel (OC)$}
        {~emphase#\ShortCascade@{$(CH) \perp (AB)$}
                      {$(OB) \perp (AB)$}}
        {or\enskip $(CH) \parallel (OB)$}
donc $(OBHC)$ est un parallélogramme
\end{Verbatim}


\bigskip
\section{L'option t}

Si on utilise l'option |t| sur la commande |\Cascade| englobante, une structure
complète de commandes |\Cascade| est alignée sur la ligne du haut.

Quand cette clé |t| est utilisée, si on souhaite ajouter du texte après la
structure, on doit le placer entre chevrons pour que ce texte soit aligné avec
la dernière accolade.


\begin{Verbatim}
\begin{enumerate}
\item \Cascade~emphase#[t]@{\ShortCascade{$(BH) \perp (AC)$}{$(OC) \perp (AC)$}}
        {or\enskip $(BH) \parallel (OC)$}
        {\Cascade{}{$(CH) \perp (AB)$}{}{$(OB) \perp (AB)$}}
        {or\enskip $(CH) \parallel (OB)$}
        ~emphase#<donc $(OBHC)$ est un parallélogramme>@
\end{enumerate}
\end{Verbatim}


\begin{enumerate}
\item \Cascade[t]{\ShortCascade{$(BH) \perp (AC)$}{$(OC) \perp (AC)$}}
        {or\enskip $(BH) \parallel (OC)$}
        {\Cascade{}{$(CH) \perp (AB)$}{}{$(OB) \perp (AB)$}}
        {or\enskip $(CH) \parallel (OB)$}
        <donc $(OBHC)$ est un parallélogramme>
\end{enumerate}



\bigskip
\section{Autres options}

\begin{itemize}
\item L'option |space-between| est une dimension LaTeX décrite sur la figure
suivante. Sa valeur initiale est de $0.5$~em. Elle s'applique à la commande
courante mais aussi à toutes les éventuelles commandes imbriquées dedans.


\item L'option |interline| peut être utilisée pour \emph{augmenter}
l'«interligne»
comme illustré sur la figure suivante. La valeur initiale de |interline| est de
$0$~pt et s'applique seulement à la commande |\Cascade| courante.

\item L'option |interline-all| change la valeur initiale de |interline| utilisée
par la commande |\Cascade| courante mais aussi toutes les éventuelles commandes
|\Cascade| imbriquées.
\end{itemize}

\vspace{1cm}

\begin{center}
\Cascade{\fbox{\hbox to 2cm{\vrule height 18 pt depth 5 pt width 0pt\hfil \texttt{\#1}\hfil
                            \tikz[remember picture] \coordinate (one) at (0,0) ;}}}%
{\fbox{\vtop{\hbox to 3cm{\tikz[remember picture] \coordinate (two) at (0,0) ;
                          \vrule height 10 pt depth 12 pt width 0pt\hfil \texttt{\#2}\hfil}
             \hbox{\tikz[remember picture] \coordinate (three) at (0,0) ;}}}}%
{\fbox{\vbox{\hbox{\tikz[remember picture] \coordinate (four) at (0,0) ;}%
             \hbox to 2.5cm{\vrule height 15 pt depth 14 pt width 0pt\hfil \texttt{\#3}\hfil}}}}%
{\fbox{\hbox to 1.7cm{\vrule height 12 pt depth 5 pt width 0pt\hfil \texttt{\#4}\hfil}}}
\begin{tikzpicture}[remember picture,overlay,blue]
\draw (one) to ++(0,1cm) coordinate (A) to ++(0,2mm) ;
\draw (two) to ++(0,1cm) coordinate (B) to ++(0,2mm) ;
\draw (A) -- (B) ;
\draw[<-] (A) -- ++(-1mm,0) ;
\draw[<-] (B) -- ++(1mm,0) ;
\path (B) node [right=2mm] {\texttt{space-between}} ;
\draw (three) to ++(-3cm,0) coordinate (C) to ++(-2mm,0) ;
\draw (four) -| coordinate (D) (C) ; 
\draw (D) to ++(-2mm,0) ;
\draw[<-] (C) -- ++(0,1mm) ;
\draw[<-] (D) -- ++(0,-1mm) ; 
\path ($(C)!0.5!(D)$) node [left=2mm] {l'«interligne»} ;
\end{tikzpicture}
\end{center}

\begin{Verbatim}
\Cascade[~emphase#interline=4mm@]{\ShortCascade{A}{B}}{E}{\ShortCascade{C}{D}}{F} G
\end{Verbatim}

\begin{center}
\Cascade[interline=4mm]{\ShortCascade{A}{B}}{E}{\ShortCascade{C}{D}}{F} G
\end{center}

\bigskip
\begin{Verbatim}
\Cascade[~emphase#interline-all=4mm@]{\ShortCascade{A}{B}}{E}{\ShortCascade{C}{D}}{F} G
\end{Verbatim}

\begin{center}
\Cascade[interline-all=4mm]{\ShortCascade{A}{B}}{E}{\ShortCascade{C}{D}}{F} G
\end{center}

\interitem
Ces options peuvent être fixées au niveau du document avec la commande 
|\CascadeOptions|. Dans ce cas, la portée de ces déclarations est le groupe TeX
courant (ces déclarations sont, comme on dit parfois «semi-globales»).

\bigskip
\section{La commande \textbackslash Edacsac}

La commande |\Edacsac| (\textsl{cascade} écrit à l'envers) est similaire à la commande |\Cascade| mais avec des
accolades dans l'autre sens. La clé~|t| n'est pas disponible dans cette commande. 

\begin{Verbatim}
Singularity 
~emphase#\Edacsac@
  {elementary}
    {
      ~emphase#\Edacsac@
        {non-degenerate elementary}
          {~emphase#\ShortEdacsac@{hyperbolic}{non-hyperbolic}}
        {degenerate elementary}
          {}
      }
  {non-elementary}
    {~emphase#\ShortEdacsac@{Nilpotent}{Higher order}}
\end{Verbatim}


Singularity 
\Edacsac
  {elementary}
  {
    \Edacsac
      {non-degenerate elementary}
      {\ShortEdacsac{hyperbolic}{non-hyperbolic}}
      {degenerate elementary}
      {}
    }
  {non-elementary}
  {\ShortEdacsac{nilpotent}{higher order}}


\bigskip
\section{Remarque technique}

Les commandes de \pkg{cascade} sont conçues pour fournir par défaut des résultats similaires à ceux fournis par les
environnements de l'\pkg{amsmath} (et de \pkg{mathtools}), en particulier |{aligned}|.


\bigskip
\begin{BVerbatim}[baseline=c,boxwidth=10cm]
\[\left.\begin{aligned}
& A = \sqrt{a^2+b^2} \\
& B = \frac{ax+b}{cx+d}
\end{aligned}\right\}\]
\end{BVerbatim}
$\left.\begin{aligned}
& A = \sqrt{a^2+b^2} \\
& B = \frac{ax+b}{cx+d}
\end{aligned}\right\}$

\bigskip
\begin{BVerbatim}[baseline=c,boxwidth=10cm]
\ShortCascade{$\displaystyle A = \sqrt{a^2+b^2}$}
             {$B = \dfrac{ax+b}{cx+d}$}
\end{BVerbatim}
\ShortCascade{$\displaystyle A = \sqrt{a^2+b^2}$}{$B = \dfrac{ax+b}{cx+d}$}

\interitem
L'extension \pkg{cascade} construit les accolades avec la paire habituelle
|\left|-|\right| de TeX. Néanmoins, les délimiteurs extensibles, en TeX, ne
peuvent pas prendre toutes les tailles. On donne, dans l'exemple suivant, les
accolades obtenues en entourant des filets verticaux de $6$~mm à $17$~mm (le code
utilise la couche programmative de LaTeX3).

\begin{Verbatim}
\int_step_inline:nnnn 6 1 {17} { $\left.\hbox{\vrule height ~#1 mm}\right\}$\quad }
\end{Verbatim}

\medskip
\begin{center}
\ExplSyntaxOn
\int_step_inline:nnnn 6 1 {17} { $\left.\hbox{\vrule height #1 mm}\right\}$\quad}
\ExplSyntaxOff
\end{center}


\section*{Autre documentation}

Le document |cascade.pdf| (fourni avec l'extension \pkg{cascade}) contient une traduction anglaise de la
documentation ici présente, ainsi que le code source commenté et un historique des versions.


\end{document}





